\subsection{An integral reduction for the other idempotent six-element source}

\begin{theorem}[Proved]
\label{CP2:A-six-constant-idempotent-integral-reduction}
Let $M$ be the six-element monoid
\[
 M=\{1,a,e_{00},e_{01},e_{10},e_{11}\},\qquad
 e_{ij}e_{kl}=e_{il},\qquad
 ae_{ij}=e_{0j},\qquad e_{ij}a=e_{i0},\qquad a^2=a.
\]
Let $S$ be any finite nonempty left-$N$, right-$M$ biset, with no
restriction on the finite monoid $N$ or on the number of states.
There is an actual quotient functor from its two-object collage $C$
to a collage $D$ with strictly fewer arrows which induces an integral
homology equivalence and a surjection on fundamental groups.
If $BC$ is simply connected, this is a homotopy equivalence.
In particular no such collage with at most eighteen arrows has
homotopy type $\mathbb{CP}^2$.  The opposite assertion also holds.
\end{theorem}
\begin{proof}
Put $e=e_{00}$ and $P=\{1,e_{00},e_{01},e_{10},e_{11}\}$.
The map
\[
 \pi:M\longrightarrow P,\qquad \pi(a)=e,
 \qquad \pi|_P=\operatorname{id}_P
\]
is a unital monoid retraction.  On $S$ write $A(s)=sa$ and $E(s)=se$.
These are idempotents with
\[
 A E=E A=E,
\]
and they commute with the left $N$ action.  The sets $ES\subseteq AS$
are nonempty.  Here $AS$ denotes the image of the function $A$, not
a left monoid action.

Form a quotient set $T$ as follows: for each $t\in AS$, identify $t$
with $Et$, and make no identifications among the states outside $AS$.
Its classes have the canonical representatives
\[
 (S\setminus AS)\ \amalg\ ES.
\]
Indeed $E$ fixes $ES$ pointwise, and its fibers inside $AS$ are exactly
the nontrivial classes of this equivalence relation.  This quotient
is left-$N$ equivariant, since $N$ commutes with $A$ and $E$.
It also carries an actual right $P$ action.  For a band element $k$
and $t=As\in AS$ one has
\[
 tk=sak=sek=(Et)k,
\]
because $ak=ek$ and $EA=E$.  Thus every defining identification is
preserved by $P$.  Moreover the classes of $As$ and $Es$ are equal
for every $s$, so the quotient map $q_S:S\to T$ is compatible with
$\pi$ on the entire source monoid $M$.  It therefore defines an
actual functor $q:C\to D$, where $D$ has corners $P,N$ and connecting
biset $T$.  This functor is the identity on the two objects and is
surjective on every hom-set.  Its target has at most
$5+|N|+|S|=|\operatorname{Mor}C|-1$ arrows.

The homological assertion is integral.  In $R=\mathbb ZM$ put
\[
 h=a-e,\qquad p=1-h=1-a+e.
\]
Direct multiplication gives
\[
 h^2=h,\qquad h\,\mathbb Z\{e_{ij}\}
   =\mathbb Z\{e_{ij}\}\,h=0,\qquad ha=ah=h.
\]
Consequently $h$ and $p$ are complementary central idempotents,
and
\[
 R=Rp\times Rh,\qquad Rp\cong\mathbb ZP,\qquad Rh=\mathbb Zh.
\]
For the middle isomorphism, $Rp$ has the integral basis $p,e_{ij}$;
the first element is its identity and the other four have exactly
the rectangular-band multiplication.  The ring map induced by
$\pi$ restricts to this isomorphism and kills $Rh$.  The trivial
left $R$ module $\mathbb Z$ lies entirely in the $p$ factor, since
the augmentations of $h,p$ are respectively $0,1$.

Let $L=\mathbb Z[S]$.  Its right summand $Lp$ has the integral basis
\[
 [s]-[As]+[Es]\quad(s\in S\setminus AS),
 \qquad [c]\quad(c\in ES).
\]
To verify spanning, the projection of a state $t\in AS$ is
$[t]p=[Et]$, while the projection of any other state is the first
displayed vector.  For independence, each such vector has coefficient
$1$ on its own state outside $AS$, and all its other terms lie inside
$AS$.  The remaining vectors are distinct state vectors in $ES$.
Under the actual quotient map $\mathbb Z[q_S]$, these vectors map
respectively to the basis elements $[s]$ and $[c]$ of $\mathbb Z[T]$.
Hence
\[
 Lp\longrightarrow\mathbb Z[T]
\]
is an integral isomorphism of left-$\mathbb ZN$, right-$\mathbb ZP$
modules, using $Rp\cong\mathbb ZP$ on the right.  No averaging or
division by a group order occurs.

The central direct-product decomposition gives the natural derived
comparison
\[
 \mathbb Z[S]\otimes_R^{\mathbf L}\mathbb Z
 \simeq
 Lp\otimes_{Rp}^{\mathbf L}\mathbb Z
 \xrightarrow{\ \cong\ }
 \mathbb Z[T]\otimes_{\mathbb ZP}^{\mathbf L}\mathbb Z.
\]
It is the comparison induced by $q_S$ and $\pi$, and is
$\mathbb ZN$ linear.  For example, project any projective left
$R$ resolution of the trivial module to the two central summands:
its $p$ part resolves $\mathbb Z$ over $Rp$, and its $h$ part is an
acyclic bounded-below projective complex and hence contractible.
The same decomposition with the trivial right module proves that
$BM\to BP$ is an integral homology equivalence.

The two-sided bridge is computed by applying
$\mathbb Z\otimes_{\mathbb ZN}^{\mathbf L}-$ to the displayed
right-action comparison.  It follows that the actual map
\[
 B(N,S,M)\longrightarrow B(N,T,P)
\]
is an integral homology equivalence as well.  The other corner map
is the identity on $BN$.  The actual two-block mapping-cylinder
description and its integral Mayer--Vietoris sequence therefore show
that $Bq:BC\to BD$ induces isomorphisms on integral homology in every
degree.

Finally, a functor which is the identity on objects and surjective
on all hom-sets is surjective on the fundamental groupoid: every
edge and reversed edge in a nerve edge-path loop has a lift with
the same endpoints.  Thus $\pi_1(BC)\to\pi_1(BD)$ is surjective.
If $BC$ is simply connected, $BD$ is too.  The integral homology
equivalence between these simply connected CW spaces is then a
homotopy equivalence.  In particular a projective-plane model with
at most eighteen arrows would give one with at most seventeen,
contrary to Theorem~\ref{CP2:A-global-lower-eighteen}.  Taking
opposite categories proves the final assertion.
\end{proof}

\begin{remark*}
This is source $1$ in the frozen five-table positive order-six
catalog.  The integral basis argument allows any number of regular
band core orbits.  It therefore removes that entire source class
from the $(6,4,8)$ budget, rather than only its actions on at most
four states.  The source with central idempotent extra actor is
separately removed by Theorem~\ref{CP2:A-extreme-idempotent-corner-reduction}.
\end{remark*}
